It is apparent from the super $T$-path formulation in Theorem~\ref{thm:main} that these $\lambda$-lengths satisfy something analogous
to the Laurent phenomenon exhibited by ordinary cluster algebras. This is summarized in the following corollary.

\begin{Corollary}\label{cor:laurent}
 Let $\tilde{\theta}_i := {\rm wt}(\sigma_i) = \sqrt{\frac{x_j}{x_kx_l}} \theta_i$ $($see Definition {\rm \ref{def:weight})}.
 For any pair of vertices~$a$,~$b$ of the polygon,
 \begin{itemize}\itemsep=0pt
 \item[$(a)$] $\lambda_{ab} \in {\mathbb R}\bigl[x_1^{\pm 1}, \dots, x_{2n+3}^{\pm 1} \,|\, \tilde{\theta}_1, \dots, \tilde{\theta}_{n+1}\bigr]$.
 In other words, each term of $\lambda_{ab}$ is the product of a~Laurent monomial
 in the $x_i$'s and a monomial in the $\tilde{\theta}_i$'s.
 \item[$(b)$] $\lambda_{ab} \in {\mathbb R}\bigl[x_{1}^{\pm \frac{1}{2}}, \dots, x_{2n+3}^{\pm \frac{1}{2}} \,|\, \theta_1, \dots, \theta_{n+1}\bigr]$.
 In other words, each term of~$\lambda_{ab}$ is the product of a~Laurent monomial in the square roots of the~$x_i$'s
 and a monomial in the~$\theta_i$'s.
 \end{itemize}
\end{Corollary}
